\chapter{Conclusion}
\label{chap:conclusion}

% ────────────────────────────────────────────────────────────
\section{Doctoral Contribution Restated}
\label{sec:contribution-restated}

This thesis set out to answer a question that is both practical and
methodological: can a research engineer build, from scratch, a
reproducible distributed platform for protein functional annotation that
is honest about the difficulty of the problem and principled in how it
measures progress?

The answer embedded in the four years of work documented here is yes,
but with a caveat that the evaluation chapter makes precise. Honest
measurement requires rigorous temporal isolation. The anc2vec feature
leakage incident, documented in ADR~D30 and Appendix~B, demonstrated
that a reranker trained without strict temporal holdout can appear to
raise NK-BPO $F_{\max}$ from 0.431 to 0.599 over the embedding-only
baseline, a gain that vanishes almost entirely once the training window
is properly quarantined from the test window. The leakage-free reranker
(v3, full 52-feature schema) instead achieves a genuine and reproducible
improvement from a baseline of 0.412 to 0.431~$F_{\max}$ in NK-BPO,
with the largest verified gains in MFO (+0.049 in LK-MFO) and CCO
(+0.024 in NK-CCO). Smaller but consistent. More important.

The doctoral contribution is threefold.

\begin{description}
  \item[PROTEA as a reproducible, operationally auditable platform.]
    An open-source distributed system implementing the full annotation
    pipeline from FASTA upload to structured \gls{go} predictions, backed
    by a versioned relational data model that links every prediction to the
    exact ontology release, annotation set, embedding configuration,
    dataset, reranker model and code commit that produced it. The platform
    is horizontally scalable through queue-based worker pools, supports
    eight \gls{plm} backends through a contracts-first plugin architecture,
    and has been deployed across development, cloud, HPC (Apptainer) and
    airgapped environments without code branching.

  \item[A leakage-free reranker pipeline that quantifies the true lift
    of feature engineering over PLM embeddings.]
    A LightGBM selective learning-to-rank reranker trained per category
    (NK, LK, PK) over a 52-feature schema combining embedding distance,
    Needleman--Wunsch and Smith--Waterman alignment statistics,
    NCBI-taxonomy distance metrics, GO-graph ancestry embeddings, and
    per-PLM PCA projections. Training follows a temporal holdout protocol
    spanning twelve consecutive GOA splits, making it reproducible and
    resistant to the leakage pathologies that inflate reported performance
    in the literature. Feature importance analysis confirms that alignment
    and taxonomy features are genuinely complementary to embedding
    distance, and not merely correlated noise.

  \item[The LAFA~v2 method submission as an independently verifiable
    performance reference.]
    Three pre-built containers were submitted to the LAFA public
    benchmark, demonstrating that PROTEA's canonical pipeline can be
    deployed and evaluated in a blinded external setting without ad-hoc
    adaptation. The experimental work also contributed to a team placement
    of second in CAFA~6, validating the embedding-based annotation
    transfer approach in an international blinded benchmark.
\end{description}

% ────────────────────────────────────────────────────────────
\section{Summary of Empirical Results}
\label{sec:empirical-summary}

\Cref{chap:evaluation} reports seven progressive experiments on the
GOA~220$\to$229 evaluation set (20,281 proteins; NK: 2,831 / LK: 3,410 /
PK: 15,313) using IA-weighted $F_{\max}$ as the primary metric.

The embedding-only baseline ($k{=}5$, aspect-separated KNN over 527,000
Swiss-Prot reference sequences embedded with ESMC~300M) reaches $F_{\max}$
values of 0.412 (NK-BPO), 0.590 (NK-MFO) and 0.668 (NK-CCO) without any
sequence alignment. This alone exceeds eggNOG-mapper v2.1.13 in 8 of 9
category--aspect cells, with gains as large as +0.306~$F_{\max}$ (NK-CCO),
establishing that \gls{plm} embeddings provide a stronger retrieval signal
than DIAMOND-based orthology transfer under this evaluation protocol.

The reranker development demonstrated that feature engineering remains
valuable even when powerful embeddings are available, but only when the
training data is correctly constructed. The v2 reranker (where alignment
and taxonomy features were inadvertently NULL during training) failed to
surpass the simple heuristic that combined those features linearly. The
v3 reranker, trained with the full 52-feature schema, achieves the best
$F_{\max}$ in 7 of 9 cells. The largest absolute gains over the
embedding-only baseline are +0.049 (LK-MFO) and +0.030 (NK-MFO), with
consistent improvements across BPO and CCO as well (\cref{tab:improvement}).

Feature importance analysis of the v3 models (\cref{tab:feature-importance})
reveals three structural findings: (i)~aggregation features
(\texttt{ref\_annotation\_density}, \texttt{go\_term\_frequency}) dominate
NK and LK; (ii)~alignment features (\texttt{identity\_nw},
\texttt{similarity\_nw}) rank 4th--6th across all three models; (iii)~raw
embedding distance falls outside the top-10 features in NK/LK, indicating
that the reranker uses embedding similarity primarily as a candidate-generation
filter rather than as the final ranking signal. This finding was not
anticipated at the start of the project and has design implications for
future feature schemas.

The leakage incident documented in ADR~D30 deserves emphasis as an empirical
finding in its own right. The pre-leakfix reranker appeared to raise
NK-BPO from 0.431 to 0.599, a gain of +0.168. After removing the anc2vec
feature leakage, the verified gain is +0.019 in the same cell. The
difference between these two numbers, 0.149~$F_{\max}$ in a single cell,
illustrates concretely why temporal holdout discipline matters and why
the PROTEA data model, which makes every experiment fully reproducible,
is not a software engineering convenience but a scientific requirement.

% ────────────────────────────────────────────────────────────
\section{Answers to the Research Questions}
\label{sec:rq-answers}

\textbf{RQ1}: \emph{Can embedding-based KNN annotation transfer match or
exceed alignment-based methods for GO term prediction?}

Yes. The embedding-only baseline surpasses eggNOG-mapper in 8 of 9
cells on the GOA~220$\to$229 evaluation, and the reranker v3 surpasses
it in all 9. The margins are largest in NK and LK, where eggNOG-mapper
suffers from a 14.5\% coverage gap and uniform confidence scores. Even
with caveats around evaluation protocol differences
(\cref{subsec:exp7-eggnogmapper}), the advantage of the embedding-based
approach at low sequence identity is clear.

\textbf{RQ2}: \emph{Do pairwise alignment statistics and taxonomic distance
provide complementary signals?}

Yes. The progression from v2 to v3 provides controlled evidence: holding
the model architecture, hyperparameters and training splits fixed, making
alignment and taxonomy features available during training improved
$F_{\max}$ in 8 of 9 cells, with the largest gain in LK-MFO (+0.032).
The v2 model, without those features, failed to surpass the simple heuristic
that combined them linearly; the v3 model, with the same features, exceeded
it. The 52-feature schema is the mechanism by which this complementarity
is captured.

\textbf{RQ3}: \emph{How can a reproducible, scalable annotation pipeline
be designed?}

Through a versioned relational data model in which every prediction is
permanently linked to an immutable \texttt{EmbeddingConfig} UUID, a
versioned \texttt{AnnotationSet} (GOA release number), and a dated
\texttt{OntologySnapshot} (OBO release date). The temporal holdout
training pipeline is fully parameterised: re-running the same payload
against the same database state reproduces identical models and results.
The leakage incident demonstrated that this reproducibility is load-bearing,
not decorative: it was the ability to re-run all experiments cleanly
after the leakage fix that made the corrected numbers trustworthy.

\textbf{RQ4}: \emph{How can the journey of a research platform, with its
dead ends and pivots, be captured operationally?}

Through the \texttt{Job.findings} and \texttt{ExperimentRun} layers and
their associated free-form comments, which provide an auditable trace from
which the evaluation chapter was distilled. ADR~D30 (operational insights
appendix) formalises this: consequential incidents such as the schema SHA
drift, the anc2vec leakage, and the cafaeval coverage bug are documented
in Appendix~B rather than buried in commit history, making the platform's
empirical narrative verifiable independently of the code.

% ────────────────────────────────────────────────────────────
\section{Future Work}
\label{sec:future-work}

The following threads are grounded in work already partially shipped or
formally planned within the PROTEA development roadmap. They are organised
by domain rather than by priority, since priority will shift with
experimental feedback.

\subsection{F-MIL: Interpretable Predictions via Residue-Level Attention}
\label{subsec:fw-mil}

The current prediction pipeline produces GO term scores at the
protein level: a confidence value per (query protein, GO term) pair.
This is sufficient for ranking but provides no interpretability at the
residue level, a gap that limits utility for structural biologists and
drug target analysts who want to know \emph{which} regions of the sequence
drive the functional assignment.

The F-MIL thread addresses this by framing annotation as a multiple
instance learning (MIL) problem over the residue embedding sequence
produced by the \gls{plm} encoder. Each residue position constitutes
a bag instance; a permutation-invariant attention pooling layer produces
a protein-level representation that retains per-residue attention weights
as a natural interpretability signal. The protea-contracts package has
already shipped the \texttt{MILBackendContract} abstract base class and
the ESM residue-output backend (\texttt{MIL.1a}) on develop. The next
steps are: (i)~integrate residue-level predictions into the annotation
pipeline as an optional backend mode; (ii)~expose per-residue attention
weights through the REST API and the web interface; (iii)~validate that
the attention pattern identifies known active-site residues for a set
of enzymes with experimental structural data.

This thread does not require a new PLM family. It reuses the existing
ESMC~300M backend's residue-output endpoint, which is already accessible
through the contracts interface.

\subsection{F-IP: InterProScan Ensemble}
\label{subsec:fw-ip}

The evaluation in \cref{chap:evaluation} demonstrates that embedding-based
KNN transfer outperforms DIAMOND orthology (eggNOG-mapper) on the
GOA~220$\to$229 protocol. A natural next step is to incorporate
InterProScan~\cite{blum2021interpro} as a complementary annotation source
rather than a competitor. InterProScan integrates over a dozen family and
domain databases (Pfam, PRINTS, ProSite, TIGRFAMs and others) and infers
GO terms via curated mappings. Its signal is structurally different from
embedding-based KNN: it is rule-based, interpretable, and highly precise
for well-characterised protein families, but sparse for novel sequences.

The F-IP plan (IP.1--IP.4) covers four steps: (IP.1)~deploy InterProScan
as a PROTEA annotation source plugin via the existing
\texttt{protea-sources} entry point mechanism; (IP.2)~generate predictions
for the GOA~220$\to$229 test set and compute per-cell $F_{\max}$;
(IP.3)~train a reranker that combines InterProScan GO evidence with the
embedding-based features already in the 52-feature schema; (IP.4)~evaluate
the ensemble on a held-out split and report the delta relative to both the
embedding-only baseline and the eggNOG-mapper comparison.

The expected benefit is largest for PK proteins, where the protein already
has annotations in some namespaces: InterProScan family membership is a
strong prior for what additional terms are plausible.

\subsection{F-OPS: Production Hardening}
\label{subsec:fw-ops}

The current deployment covers the full pipeline but remains a
single-node system managed by \texttt{manage.sh}. Three hardening
threads are in progress.

The T-OPS bundle (\texttt{T5.1--T5.5}) covers structured logging
(JSON lines), metrics export (Prometheus), alerting (alert manager
rules for queue depth and worker error rate), distributed tracing
(OpenTelemetry spans from HTTP request to queue operation to database
write), and a health dashboard. The observability gap is the most
operationally significant limitation: without structured logs, diagnosing
worker failures or queue backlogs requires shell access to the deployment
host.

Authentication (\texttt{T5.6}) adds a token-based auth layer to the
REST API. The current API is unauthenticated and suitable only for
single-tenant or firewalled deployments. Token auth is a prerequisite
for any multi-tenant or public-facing instance, and is also required
before PROTEA can be offered as a web service to collaborators at
external institutions.

The airgapped Apptainer bundle (ADR~D27) is deferred to post-defense.
The priority is to ship T-OPS and T5.6 so that the platform is
operationally defensible in a production context, not merely a
research prototype that requires developer attention to keep running.

\subsection{LAFA v2 Follow-Ups}
\label{subsec:fw-lafa}

Three containers were submitted to the LAFA public benchmark (F-LAFA.1
through F-LAFA.3 as described in \cref{chap:introduction}). The
leakage-free reranker was not yet available at the time of the initial
submissions: those containers used the pre-leakfix feature schema. The
F-LAFA.3 thread covers rebuilding the LAFA submission container with the
corrected v3 reranker and the full 52-feature schema, so that the LAFA
leaderboard result reflects the same pipeline as the evaluation chapter.

F-LAFA.4 will run a full ablation study on the LAFA evaluation set,
paralleling the GOA~220$\to$229 experiments in \cref{chap:evaluation}
but using the LAFA-held-out proteins as the test set. This provides an
independent verification of the reranker's generalisation beyond the
single GOA split used in this thesis. It also tests whether the feature
importance hierarchy observed in \cref{tab:feature-importance} holds on
a different protein set.

% ────────────────────────────────────────────────────────────
\section{Engineering Wrap-Up}
\label{sec:engineering-wrapup}

Three architectural decision records close the engineering narrative.

\textbf{ADR~D29 (Release pipeline).} The per-repo SemVer release pipeline
driven by Conventional Commits and \texttt{release-please} automates
version tagging, changelog generation, and container image publication
to GitHub Container Registry. The release pipeline is a prerequisite for
the T-OPS observability work (F-OPS) because structured versioning is
required before cross-repo integration tests can be reliably anchored
to a stable tag.

\textbf{ADR~D30 (Operational insights appendix).} This ADR formalises
the decision to document consequential incidents (schema SHA drift,
anc2vec feature leakage, cafaeval coverage bug) in Appendix~B rather
than in commit messages or design documents. The rationale is auditability:
incidents that affected the empirical results of the thesis should be
in a location that a reader can find without access to the git log.
The leakage incident described in \cref{sec:empirical-summary} is the
primary example. ADR~D30 makes the operational narrative an explicit part
of the thesis artefact rather than an implicit side effect of version
control.

\textbf{ADR~D31 (Method Object reframe).} The refactor strategy for large
batch operation classes, extracting cohesive sub-clusters as
\texttt{\_<Cluster>Runner} helper classes rather than wrapping
already-budgeted \texttt{execute()} methods, is codified in D31. This
decision directly shapes how the F-MIL and F-IP operation classes will
be structured when they land in \texttt{protea-runners}: new operations
follow the same sub-cluster pattern, keeping the executor graph navigable
as the operation catalogue grows from its current 17-router topology.

% ────────────────────────────────────────────────────────────
\section{Closing Remarks}
\label{sec:closing}

PROTEA began as a question about reproducibility in computational biology
and ended as a demonstration that the question is load-bearing. The
leakage incident is the most direct evidence: a platform without a
versioned data model and reproducible training pipelines cannot distinguish
a genuine +0.049~$F_{\max}$ gain from an artefact of data contamination.
A platform with those properties can.

The functional annotation gap is not close to being solved. Fewer than
0.02\,\% of deposited proteins carry experimentally validated annotations,
and that fraction shrinks as sequencing throughput grows. PROTEA's
contribution is not to close the gap but to provide a foundation from
which incremental, verifiable progress can be made, and to demonstrate
on a concrete evaluation set that embedding-based annotation transfer
with selective feature-engineered reranking is a meaningful step in
the right direction.

The platform is released as open-source software at
\url{https://github.com/frapercan/protea}. Francisco Miguel Pérez Canales
is its author and sole maintainer.
